\section{Referencial Teórico}

Esta seção apresenta os conceitos que sustentam o trabalho.

\subsection{Dados abertos governamentais}

Dados abertos governamentais são informações produzidas ou mantidas pelo poder público e disponibilizadas para que outras pessoas possam usá-las de novo.
O objetivo costuma ser aumentar a transparência e permitir que esses dados sejam reaproveitados.
\citeonline{molodtsov2024} estudaram a usabilidade de portais de dados abertos em 33 países.
O quadro que eles aplicam não olha só se o arquivo está disponível para \textit{download}.
Também entra se a pessoa consegue explorar o dado e entendê-lo, mesmo sem prática com planilha.
\citeonline{ansari2022}, ao revisar pesquisas sobre visualização de dados governamentais, observam um ponto parecido: publicar o arquivo não faz, por si só, com que o dado seja usado.
Muita gente não consegue ler o formato, cruzar bases diferentes e interpretar a tabela.
Uma aplicação que junta as bases e mostra mapa, gráfico e filtro pode tornar essa consulta mais simples.
Esse é o caminho deste trabalho: caracterizar as empresas de TI de Contagem a partir de bases abertas nacionais, e não a partir de um arquivo municipal já pronto.

\subsection{System Usability Scale (SUS)}

Usabilidade, neste trabalho, é a facilidade de usar a interface.
A pessoa encontra a informação, entende o que está vendo e consegue filtrar sem se perder.
A System Usability Scale (SUS), proposta por \citeonline{brooke1996}, é um questionário de dez itens em escala Likert de cinco pontos, com afirmações alternadas entre positivas e negativas.
Após o uso do sistema, cada participante responde aos itens; a pontuação individual é convertida para uma escala de 0 a 100.
Essa medida única resume a percepção global de usabilidade e permite comparar o resultado com faixas de interpretação usadas na literatura.
Este trabalho aplica o SUS à aplicação desenvolvida, após os participantes percorrerem o mapa, os filtros, os indicadores e a tabela.

\subsection{Classificação Nacional de Atividades Econômicas (CNAE)}

A identificação de um setor em cadastros públicos depende de uma classificação padronizada das atividades.
No Brasil, a Classificação Nacional de Atividades Econômicas (CNAE), mantida pelo IBGE, é o instrumento usado pelo sistema estatístico e pelos órgãos que registram empresas \cite{ibge2020}.
A CNAE organiza as atividades em níveis, da seção até a subclasse.
No CNPJ, cada estabelecimento tem um código de atividade principal.

Neste trabalho, empresa de TI é o estabelecimento cuja atividade principal está na divisão 62, de serviços de tecnologia da informação, na versão CNAE-Subclasses 2.3.
Os códigos incluídos são:
\begin{itemize}
    \item 6201-5/01, desenvolvimento de programas sob encomenda;
    \item 6201-5/02, web design;
    \item 6202-3/00, programas customizáveis;
    \item 6203-1/00, programas não customizáveis;
    \item 6204-0/00, consultoria em tecnologia da informação;
    \item 6209-1/00, suporte técnico e outros serviços em TI.
\end{itemize}
Não entram a divisão 63 (hospedagem, tratamento de dados e portais), o comércio e a fabricação de equipamentos, o conserto de computador nem o curso de informática.
A regra usa só a atividade principal, para a lista não crescer com códigos secundários de empresas de outro ramo.
